\providecommand{\hardq}{\refstepcounter{enumi}\item[\labelenumi$^{\dagger}$]}
\providecommand{\hardp}{\refstepcounter{enumii}\item[\labelenumii$^{\dagger}$]}
\noindent Questions, or parts of questions, marked with a dagger ($\dagger$) are more challenging than the rest. They can be attempted, or their solutions simply read through; everything unmarked should be attempted.

\begin{enumerate}

\item In Lecture~18 it was shown that, for a P-wave ray in a spherically symmetric earth model of radius $b$ whose turning 
point lies at radius $r_{t}$, the epicentral angle and travel time are
\begin{equation*}
    \Delta(q)=\int_{r_{t}}^{b}\frac{2\alpha q}{r\sqrt{r^{2}-\alpha^{2}q^{2}}}\dd r, \quad 
    T(q)=\int_{r_{t}}^{b}\frac{2r}{\alpha\sqrt{r^{2}-\alpha^{2}q^{2}}}\dd r,
\end{equation*}
where $q=r\sin\theta/\alpha$ is the ray parameter, $\theta$ being the angle between the ray and the radial direction, and $r_{t}$ satisfies $r_{t}/\alpha(r_{t})=q$. This question 
shows how such a travel-time curve can be inverted for the velocity structure, a method due to Herglotz and Wiechert. Throughout, 
define $\eta(r)=r/\alpha(r)$ and assume that $\eta$ increases monotonically with $r$.
\begin{enumerate}
\item Show that $\tau(q)=T(q)-q\Delta(q)$ can be written
\begin{equation*}
    \tau(q)=2\int_{r_{t}}^{b}\sqrt{\eta^{2}-q^{2}}\,\frac{\dd r}{r},
\end{equation*}
and hence that $\dd\tau/\dd q=-\Delta$. [Remember that the lower limit $r_{t}$ depends on $q$, so that Leibniz's rule for differentiating an integral with respect to a parameter is needed; consider the value of the integrand at the turning point.] Deduce that $\dd T/\dd\Delta=q$ along the travel-time curve, so that the ray parameter of 
an arrival can be measured from the local slope of the observed curve.
\item Using $\eta$ as the integration variable, show that
\begin{equation*}
    \Delta(q)=2\int_{q}^{\eta_{b}}\frac{q}{\sqrt{\eta^{2}-q^{2}}}\frac{\dd\ln r}{\dd\eta}\dd\eta,
\end{equation*}
where $\eta_{b}=\eta(b)$.
\hardp Let $\eta_{1}$ be the ray parameter of the ray turning at some radius $r_{1}$. Multiply the result of (b) by 
$(q^{2}-\eta_{1}^{2})^{-1/2}$, integrate over $q$ from $\eta_{1}$ to $\eta_{b}$, and exchange the order of integration to show that
\begin{equation*}
    \ln\frac{b}{r_{1}}=\frac{1}{\pi}\int_{\eta_{1}}^{\eta_{b}}\frac{\Delta(q)}{\sqrt{q^{2}-\eta_{1}^{2}}}\dd q.
\end{equation*}
You may use the identity
\begin{equation*}
    \int_{\eta_{1}}^{\eta}\frac{q\,\dd q}{\sqrt{(\eta^{2}-q^{2})(q^{2}-\eta_{1}^{2})}}=\frac{\pi}{2}.
\end{equation*}
By integrating by parts, show that this can equally be written in terms of the observed travel-time curve as
\begin{equation*}
    \ln\frac{b}{r_{1}}=\frac{1}{\pi}\int_{0}^{\Delta_{1}}\cosh^{-1}\!\left[\frac{q(\Delta)}{\eta_{1}}\right]\dd\Delta,
\end{equation*}
where $\Delta_{1}=\Delta(\eta_{1})$ and $q(\Delta)=\dd T/\dd\Delta$.
\item Explain how these results allow $\alpha(r)$ to be determined from a travel-time curve, working downwards from the surface. 
What goes wrong if there is a region in which the velocity decreases with depth so rapidly that $\eta$ is no longer monotonic? 
And what is seen in the travel-time curve at a discontinuity across which the velocity increases with depth?
\end{enumerate}

\item Consider the regularised least squares solution to a linear(ised) inverse problem obtained by minimising the function
\begin{equation*}
    J(\mathbf{m})=\frac{1}{2}(\mathbf{A}\mathbf{m}-\mathbf{d})^{T}\mathbf{C}^{-1}(\mathbf{A}\mathbf{m}-\mathbf{d})+\frac{\lambda}{2}\mathbf{m}^{T}\mathbf{B}\mathbf{m}.
\end{equation*}
Suppose that the data take the form $\mathbf{d}=\mathbf{A}\mathbf{m}_{\mathrm{in}}$ with $\mathbf{m}_{\mathrm{in}}$ a specified input model. Show that the resulting least squares solution, $\mathbf{m}_{\mathrm{out}}$, takes the form
\begin{equation*}
    \mathbf{m}_{\mathrm{out}}=\mathbf{R}\mathbf{m}_{\mathrm{in}},
\end{equation*}
where $\mathbf{R}$ is a matrix to be determined. Discuss the significance of this so-called resolution matrix, explaining why, in particular, it would be desirable for its value to be close to the identity matrix.

\item

Consider a linear inverse problem of the standard form
\begin{equation*}
    \mathbf{d} = \mathbf{A}\mathbf{m}+\mathbf{e},
\end{equation*}
with $\mathbf{d}$ an $n$-dimensional data vector, $\mathbf{m}$ an $m$-dimensional model vector, $\mathbf{A}$ an $n\times m$ matrix, and $\mathbf{e}$ an $n$-dimensional vector of random errors. We assume $m > n$, which
is to say that there are more unknown model parameters than available data.
We further suppose that the data have been processed such that the components of the random errors are uncorrelated and have unit standard deviations (this corresponds to the data covariance, $\mathbf{C}$, being the $n\times n$ identity matrix).



\begin{enumerate}

    \item A model $\mathbf{m}$ is said to be compatible with the data if
\begin{equation*}
    \|\mathbf{A}\mathbf{m}-\mathbf{d}\|^{2} \le \chi^{2}_{c},
\end{equation*}
where $\chi^{2}_{c}$ is a critical value for the so-called chi-squared statistic. What type of geometric region of the model space is defined by this inequality?

\item Supposing that $\|\mathbf{d}\|^{2} > \chi^{2}_{c}$, the region of the model space defined by the inequality in (a) does not contain the origin. In such
cases, the closest point within this region to the origin must lie on the boundary set defined by
\begin{equation*}
    \|\mathbf{A}\mathbf{m}-\mathbf{d}\|^{2} = \chi^{2}_{c}.
\end{equation*}
Using the method of Lagrange multipliers, we can determine this \emph{minimum norm solution} by considering the Lagrangian
\begin{equation*}
    \mathcal{L}(\mathbf{m},\mu)=\|\mathbf{m}\|^{2}+\mu(\|\mathbf{A}\mathbf{m}-\mathbf{d}\|^{2}-\chi^{2}_{c}),
\end{equation*}
where $\mu$ is a Lagrange multiplier. Derive the equations for $\mathbf{m}$ and $\mu$ that must hold at the minimum norm solution.
Discuss the relationship between this approach and regularised least squares.
\end{enumerate}


\hardq Consider the deformation of an isotropic elastic body due to an applied surface load, $\sigma$. The relevant linearised equilibrium equations are
\begin{gather*}
    -\frac{\partial T_{ij}}{\partial x_{j}}=0, \\
    T_{ij}=\lambda\frac{\partial u_{k}}{\partial x_{k}}\delta_{ij}+\mu\left(\frac{\partial u_{i}}{\partial x_{j}}+\frac{\partial u_{j}}{\partial x_{i}}\right),
\end{gather*}
within the reference body $M$, and subject to boundary conditions
\begin{equation*}
    T_{ij}\hat{n}_{j}=\sigma g_{i}
\end{equation*}
on its boundary, $\partial M$.\footnote{For a purely traction-loaded static problem the load must exert no net force or torque and the displacement is determined only up to a rigid-body motion; we ignore this subtlety here.} Here $g_{i}$ is the acceleration due to gravity prior to the deformation. Note that self-gravitation is being ignored within these equations.

Suppose that observations, $\mathbf{u}_{r}^{\mathrm{obs}}$, of the displacement vector have been made at surface locations $\mathbf{x}_{r}$ for $r=1,\dots,N$. An inverse problem could then be formulated to estimate the Lamé parameters, $\lambda$ and $\mu$, as well as the load, $\sigma$. To this end, we can seek to minimise the misfit
\begin{equation*}
    J=\frac{1}{2}\sum_{r=1}^{N}\frac{1}{\varepsilon_{r}^{2}}\|\mathbf{u}(\mathbf{x}_{r})-\mathbf{u}_{r}^{\mathrm{obs}}\|^{2},
\end{equation*}
where $\varepsilon_{r}$ is the standard deviation for the $r$th observation. Adapt the approach within Lecture~20 to derive sensitivity kernels for $J$ with respect to $\lambda$, $\mu$, and $\sigma$, showing how they can be calculated through one solution of the forward problem and one of an appropriate adjoint problem.





\item Recall that the Lagrangian for a self-gravitating elastic body is given by
\begin{equation*}
    \mathcal{L}(\mathbf{x},t,\bphi,\mathbf{v},\mathbf{F})=\frac{1}{2}\rho(\mathbf{x})\|\mathbf{v}\|^{2}-W(\mathbf{x},t,\mathbf{F})-\frac{1}{2}\rho(\mathbf{x})\zeta(\mathbf{x},t),
\end{equation*}
where $\zeta$ is the referential gravitational potential
\begin{equation*}
    \zeta(\mathbf{x},t)=-G\int_{M}\frac{\rho(\mathbf{x}')}{\{[\varphi_{i}(\mathbf{x},t)-\varphi_{i}(\mathbf{x}',t)][\varphi_{i}(\mathbf{x},t)-\varphi_{i}(\mathbf{x}',t)]\}^{\frac{1}{2}}}\dd^{3}\mathbf{x}'.
\end{equation*}
By calculating the appropriate functional derivative, show that
\begin{equation*}
    \frac{\delta\mathcal{L}}{\delta\varphi_{i}}=\rho\gamma_{i},
\end{equation*}
where $\gamma_{i}$ is the referential gravitational acceleration.

\item Using results from Lecture~23 along with residue calculus (which you don't need to know for this course), it can be shown that the time-domain solution of the elastodynamic equations in a non-rotating earth model can be expressed as the following eigenfunction expansion
\begin{equation*}
    \mathbf{u}(\mathbf{x},t)=\sum_{km}\left\{\int_{0}^{t}\frac{\sin[\omega_{k}(t-t')]}{\omega_{k}}\int_{M}\overline{S}_{ij}(\mathbf{x}',t')\frac{\partial|km\rangle_{i}^{*}}{\partial x'_{j}}(\mathbf{x}')\dd^{3}\mathbf{x}'\dd t'\right\}|km\rangle(\mathbf{x}),
\end{equation*}
where $|km\rangle_{i}$ denotes the $i$th Cartesian component of the eigenfunction.
\begin{enumerate}
\item Suppose that the stress glut is non-zero only within a region of dimension $L$ about a point $\mathbf{x}_{s}$, as it would be for 
an earthquake on a fault of that size. By expanding $\partial|km\rangle_{i}^{*}/\partial x'_{j}$ in a Taylor series about $\mathbf{x}_{s}$, show that 
the contribution of a mode to the sum can be approximated by replacing the stress glut with the \emph{point source}
\begin{equation*}
    \overline{S}_{ij}(\mathbf{x},t)=M_{ij}(t)\delta(\mathbf{x}-\mathbf{x}_{s}), \quad M_{ij}(t)=\int_{M}\overline{S}_{ij}(\mathbf{x},t)\dd^{3}\mathbf{x},
\end{equation*}
provided that $L$ is small compared with the length scale over which the strain of the eigenfunction varies. Estimate this 
length scale in terms of the eigenfrequency and a wave speed, and hence explain why the point-source approximation is a statement 
about the frequencies retained in the seismograms. For a fault of length $100\,\mathrm{km}$ and a shear wave speed of 
$5\,\mathrm{km\,s^{-1}}$, below roughly what frequency is the approximation reasonable?
\item Simplify the eigenfunction expansion in the case of a point source with a step-function time dependence, 
\begin{equation*}
    \overline{S}_{ij}(\mathbf{x},t)=M_{ij}\delta(\mathbf{x}-\mathbf{x}_{s})H(t-t_{s}),
\end{equation*}
with $M_{ij}$ known as the moment tensor, $\mathbf{x}_{s}$ the source location, and $t_{s}$ the source time. Describe briefly and qualitatively how an inverse problem might be formulated to estimate the ten source parameters $(M_{ij},\mathbf{x}_{s},t_{s})$ from recorded seismograms.
\end{enumerate}

\item Consider the eigenvalue problem for a rotating earth model
\begin{equation*}
    -\omega^{2}\langle \mathbf{w}|P|\mathbf{s}\rangle+\ii\omega\langle \mathbf{w}|W|\mathbf{s}\rangle+\langle \mathbf{w}|H|\mathbf{s}\rangle=0,
\end{equation*}
where $\mathbf{s}$ is the eigenfunction, $\omega$ the eigenfrequency, and $\mathbf{w}$ an arbitrary test function. Assuming the eigenfunction normalisation condition
\begin{equation*}
    \langle \mathbf{s}|P|\mathbf{s}\rangle=1,
\end{equation*}
show that the eigenfrequency can be written
\begin{equation*}
    \omega=\frac{1}{2}\ii\langle \mathbf{s}|W|\mathbf{s}\rangle\pm\sqrt{\langle \mathbf{s}|H|\mathbf{s}\rangle-\frac{1}{4}\langle \mathbf{s}|W|\mathbf{s}\rangle^{2}}.
\end{equation*}
Use this result to discuss qualitatively the stability of a rotating earth model, including, in particular, the role of the Coriolis and centrifugal forces.

\item Consider a non-rotating planet of radius $a$ with uniform density $\rho$, and let $M_{\mathrm{p}}$ be its mass.
\begin{enumerate}
\item Show that the gravitational potential energy (or binding energy) of the planet is
\begin{equation*}
    E_{g}=-\frac{3GM_{\mathrm{p}}^{2}}{5a}.
\end{equation*}
\item Assuming the planet to be in hydrostatic equilibrium, show that its pressure is
\begin{equation*}
    p_{0}(r)=\frac{2\pi}{3}G\rho^{2}(a^{2}-r^{2}),
\end{equation*}
and hence verify that $3\int_{M}p_{0}\dd^{3}\mathbf{x}=-E_{g}$.
\item Suppose that the planet is compressed uniformly, so that its radius becomes $a(1-\epsilon)$ with $\epsilon$ small and positive,
the density remaining uniform. Determine the change in the gravitational potential energy correct to second order in $\epsilon$.
\item Let the material respond to compression through $p=p_{0}+\kappa s$, where $s$ is the fractional decrease in the volume of a
material element and $\kappa$ is a constant bulk modulus. Recalling that the work done in compressing a material element of initial
volume $V_{0}$ is $V_{0}\int p\,\dd s$, show that the change in the internal energy of the planet is
\begin{equation*}
    \Delta E_{\mathrm{int}}=(3\epsilon-3\epsilon^{2})\int_{M}p_{0}\dd^{3}\mathbf{x}+\frac{9}{2}\kappa V\epsilon^{2}+O(\epsilon^{3}),
\end{equation*}
where $V$ is the volume of the planet.
\item Show that the total change in energy vanishes at first order in $\epsilon$, and explain why this had to be so. Show that
the second-order change is positive, and hence the planet stable with respect to this perturbation, if and only if
\begin{equation*}
    \kappa>\frac{4}{3}\bar{p}_{0}=\frac{16\pi}{45}G\rho^{2}a^{2},
\end{equation*}
where $\bar{p}_{0}$ is the mean pressure. Express this condition in terms of the sound speed $\sqrt{\kappa/\rho}$ and the surface
gravity $g$, and evaluate it for $\rho=5500\,\mathrm{kg\,m^{-3}}$ and $a=6371\,\mathrm{km}$. Comment on the result in
relation to the bulk moduli within the earth model PREM, which range from roughly $130\,\mathrm{GPa}$ in the upper mantle to
$650\,\mathrm{GPa}$ at the base of the mantle.
\end{enumerate}

\item This question concerns the splitting of a multiplet by the Earth's rotation, using the first-order theory of Lecture~24. Ignore 
the centrifugal force and the Earth's ellipticity, so that within the multiplet the only perturbation is the Coriolis form $W$, and take 
the rotation vector to be $\bm{\Omega}=\Omega\hat{\mathbf{z}}$.
\begin{enumerate}
\item It can be shown that the Coriolis matrix elements within a multiplet of degree $l$ are diagonal in the order $m$, taking the form
\begin{equation*}
    \langle km'|W|km\rangle=-2\ii m\Omega\beta\,\delta_{mm'},
\end{equation*}
where $\beta$ is a real number depending on the radial eigenfunctions of the multiplet. Show that the singlets $|km\rangle$ are then 
eigenfunctions of the first-order splitting problem, with eigenfrequencies
\begin{equation*}
    \omega_{m}=\omega_{k}+m\Omega\beta, \quad -l\le m\le l,
\end{equation*}
so that the multiplet is split into $2l+1$ equally spaced singlets, and verify that the diagonal sum rule holds. Show that this 
is consistent with the result of question 7.
\hardp Consider a toroidal multiplet with singlets of the form
\begin{equation*}
    |km\rangle=W(r)\mathbf{C}_{lm}(\theta,\varphi)
\end{equation*}
(the radial eigenfunction $W(r)$, not to be confused with the Coriolis operator $W$), normalised so that
\begin{equation*}
    l(l+1)\int_{0}^{b}\rho W^{2}r^{2}\dd r=1.
\end{equation*}
Show that $\beta=1/[l(l+1)]$. You may use 
\begin{equation*}
    \mathbf{C}_{lm}=-\hat{\mathbf{r}}\times\nabla_{1}Y_{lm}, \quad 
    \nabla_{1}Y_{lm}=\hat{\bm{\theta}}\,\frac{\partial Y_{lm}}{\partial\theta}+\frac{\hat{\bm{\varphi}}}{\sin\theta}\frac{\partial Y_{lm}}{\partial\varphi},
\end{equation*}
the dependence $Y_{lm}\propto\ee^{\ii m\varphi}$, and the fact that $Y_{lm}$ vanishes at the poles for $m\ne0$. 
[For a spheroidal multiplet a similar calculation gives 
$\beta=\int_{0}^{b}\rho(2UV+V^{2})r^{2}\dd r$ with the normalisation $\int_{0}^{b}\rho[U^{2}+l(l+1)V^{2}]r^{2}\dd r=1$, 
which for ${}_{0}S_{2}$ in PREM gives $\beta\approx0.40$ and for ${}_{0}S_{3}$ gives $\beta\approx0.19$.]
\item The Earth rotates once per sidereal day, $86164\,\mathrm{s}$. Estimate the spacing between adjacent singlets of 
${}_{0}T_{2}$, ${}_{0}S_{2}$ and ${}_{0}S_{3}$, and compare with the observed spectra in Fig.~6 of Lecture~24. Why might the observed 
singlets not be exactly equally spaced?
\end{enumerate}

\end{enumerate}
